% Research-only source-aware proof for P06.1, 30 August 2026.
% Not included by the live manuscript.  No novelty or priority claim is made.

\begin{proof}[Proof of \Cref{lem:p061-product-open-tube}]
Let
\[
 \epsilon=\epsilon_a\widehat\otimes1:P\longrightarrow B.
\]
We first record the elementary uniform estimate used below.  If
\(G\in\ker\epsilon\) and \(\eta>0\), there is a strict affinoid
neighbourhood \(U\) of \(a\) such that
\begin{equation}
\label{eq:p061-product-kernel-small}
       \sup_{U\times_K\mathcal M(B)}|G|<\eta.
\end{equation}
Indeed, choose a finite tensor
\(T=\sum_{i=1}^r a_i\otimes b_i\) sufficiently close to \(G\), and put
\[
                  b_0=\epsilon(T)\in B.
\]
Boundedness of \(\epsilon\) and the equality \(\epsilon(G)=0\) make
\(\lVert b_0\rVert\) arbitrarily small with the approximation error.
Replacing \(T\) by
\[
 T-1\otimes b_0
   =\sum_{i=1}^r(a_i-a_i(a))\otimes b_i
\]
therefore still gives an arbitrarily close approximation.  Every
\(a_i-a_i(a)\) vanishes at
\(a\).  Nontriviality of the valuation supplies arbitrarily small strict
rational neighbourhoods on which all of them have the required small
supremum norm.  The ultrametric inequality proves
\eqref{eq:p061-product-kernel-small}.

Choose a basic ordinary neighbourhood of \(s_a(z)\), contained in
\(\Omega\), of the form
\[
 \Omega_0=
 \bigcap_{i=1}^r\{|F_i|<c_i\}
 \cap
 \bigcap_{j=1}^s\{|H_j|>d_j\},
 \qquad F_i,H_j\in P.
\]
Choose numbers
\[
 |\epsilon(F_i)(z)|<c_i'<c_i,
 \qquad
 d_j<d_j'<|\epsilon(H_j)(z)|.
\]
After shrinking to one strict affinoid neighbourhood \(V\) of \(z\), the
same strict inequalities hold throughout \(V\).  Write, for every
\(L\in\{F_i,H_j\}\),
\[
        L=1\otimes\epsilon(L)+
          \bigl(L-1\otimes\epsilon(L)\bigr).
\]
The second summand belongs to \(\ker\epsilon\).  Apply
\eqref{eq:p061-product-kernel-small} simultaneously, with error smaller
than \(c_i\) for \(F_i\) and smaller than \(d_j'\) for \(H_j\).  On the
resulting \(U\times_KV\), the upper inequalities follow from the
ultrametric triangle inequality.  In the lower inequalities the error is
strictly smaller than \(|\epsilon(H_j)|\), so
\[
                     |H_j|=|\epsilon(H_j)|>d_j'>d_j.
\]
Thus \(U\times_KV\subset\Omega_0\subset\Omega\).
\end{proof}

\begin{proof}[Proof of
\Cref{prop:p061-parameterized-toric-fixed-fibre}]
Choose \(0<|\varpi|<1\) and
\(v\in N\cap\operatorname{int}(\rho)\), and put
\(\ell(u)=\langle u,v\rangle\).  The weighted domains \(U_n\) in
\eqref{eq:p061-parameterized-weighted-domain} are strict affinoid
neighbourhoods of \(a\): for finite monoid generators
\(u_1,\ldots,u_r\) of \(S\), the corresponding monomials give a closed
immersion \(Z\hookrightarrow\mathbf A_K^r\), and \(U_n\) is the inverse
image of the strict closed polydisc
\[
             |\chi^{u_i}|\leq
             |\varpi|^{n\ell(u_i)}
             \qquad(1\leq i\leq r).
\]
After rescaling the generators, their affinoid algebras are
\(K\langle S\rangle\) with the semigroup Gauss norm.  They are nested,
and
\[
                         \bigcap_{n\geq1}U_n=\{a\}.
\]
The compactness argument for nested closed subsets shows that they are
cofinal among ordinary neighbourhoods of \(a\).  More explicitly, if an
ordinary neighbourhood \(E\) of \(a\) contained no \(U_n\), the nonempty
closed subsets \(U_n\setminus E\) of the compact space \(U_1\) would have
nonempty intersection, contradicting the displayed identity.

Choose strict affinoid neighbourhoods \(A_0\) of \(a\) in
\(Z^{\mathrm{an}}\) and \(B_0\) of \(y\) in \(Y\).  By
\Cref{lem:p061-product-open-tube}, after shrinking there are strict
affinoid neighbourhoods \(A_1\subset A_0\) and \(V\subset B_0\) such
that
\[
                         A_1\times_KV\subset\Omega.
\]
Since \(A_1\) is a neighbourhood of \(a\), choose \(n\) with
\(U_n\subset A_1\).  It remains to prove that \(h\) is constant on every
fibre of
\[
                    p_V:U_n\times_KV\longrightarrow V.
\]
We do this over the actual completed residue field of an arbitrary point,
so no density or rigid-point argument is hidden.  Fix \(b\in V\).  The
analytic fibre is
\[
 (U_n)_{\mathscr H(b)}=
 U_n\widehat\otimes_K\mathscr H(b).
\]
After the rescaling
\(Z^u=\varpi^{-n\ell(u)}\chi^u\), its affinoid algebra is the semigroup
Tate algebra
\begin{equation}
\label{eq:p061-basechanged-semigroup-tate}
 \mathscr H(b)\langle S\rangle=
 \left\{\sum_{u\in S}c_uZ^u:
 c_u\in\mathscr H(b),\ |c_u|\longrightarrow0
 \text{ off finite subsets of }S\right\}
\end{equation}
with Gauss norm \(\max_u|c_u|\).  The Gauss norm is multiplicative: after
normalizing two nonzero series to norm one, their nonzero reductions have
nonzero product in the domain
\(\widetilde{\mathscr H(b)}[S]\subset
 \widetilde{\mathscr H(b)}[M]\).  It is therefore the supremum norm, and
the canonical reduction of \eqref{eq:p061-basechanged-semigroup-tate} is
\[
                       \widetilde{\mathscr H(b)}[S].
\]
This is an integral domain.  Hence
the Gauss norm defines a point \(\eta_{n,b}\) of
\((U_n)_{\mathscr H(b)}\), and
\[
        |f(\eta_{n,b})|=\lVert f\rVert_{\mathrm G}
                       =\sup_{(U_n)_{\mathscr H(b)}}|f|
        \qquad
        \bigl(f\in\mathscr H(b)\langle S\rangle\bigr).
\]
Thus \(\{\eta_{n,b}\}\) is itself a boundary.  The nonempty minimal
Shilov boundary is consequently the singleton
\(\{\eta_{n,b}\}\).

First extend \(h\) from \(K\) to \(\mathscr H(b)\); this preserves
classical pluriharmonicity in both directions by
\cite[Proposition~7.2(6)]{GR25Classical}.  The point \(b\) induces an
\(\mathscr H(b)\)-rational point of \(V_{\mathscr H(b)}\), and the
corresponding fibre inclusion
\[
 (U_n)_{\mathscr H(b)}
 \longrightarrow (U_n\times_KV)_{\mathscr H(b)}
\]
is an \(\mathscr H(b)\)-analytic morphism.  Arbitrary pullback preserves
classical pluriharmonicity by
\cite[Proposition~7.2(7)]{GR25Classical}.  Thus the resulting pullback
\(h_b\) is pluriharmonic on the fibre.  The affinoid maximum--minimum
principle \cite[Proposition~7.3]{GR25Classical} says that both extrema of
\(h_b\) occur on its singleton Shilov boundary, so \(h_b\) is constant.
By the standard description of analytic fibres
\cite[\S2.7]{GR25Classical}, the underlying topological space of
\((U_n)_{\mathscr H(b)}\) identifies with the set-theoretic fibre
\(p_V^{-1}(b)\).  Hence every point of \(U_n\times_KV\) belongs to
exactly one of the analytic fibres just treated.  Its value under \(h\)
consequently equals the value at the toric-fixed section in that fibre.
This proves the pointwise identity
\[
              h|_{U_n\times_KV}=p_V^*(s_a^*h|_V).
\]
The function \(g=s_a^*h|_V\) is classically pluriharmonic by the same
functorial pullback result.  If pluriharmonic functions on \(Y\) are locally
\(\mathbf R\)-PL, shrink \(V\) once more and apply pullback stability of
\(\mathbf R\)-PL functions
\cite[Lemma~5.4(1)]{GJR26}.  For quasi-smooth \(Y\), the required
regularity is the already verified quasi-smooth P06.1 theorem.  This
proves every assertion.
\end{proof}

\begin{proof}[Proof of \Cref{thm:p061-split-toric-all-orbits}]
Fix \(x\in W\).  There is a unique face
\(\tau\preceq\sigma\) such that \(x\in O(\tau)^{\mathrm{an}}\).  Put
\[
 N_\tau=N\cap\operatorname{span}_{\mathbf R}(\tau).
\]
It is a saturated sublattice of \(N\), so choose a splitting
\(N=N_\tau\oplus N'\).  Choose
\[
 q\in\operatorname{relint}
 (\sigma^\vee\cap\tau^\perp)\cap M.
\]
Such an integral relative-interior point exists because the dual face is
rational polyhedral.  Since
\(\sigma\cap\operatorname{span}_{\mathbf R}(\tau)=\tau\), its zero face
on \(\sigma\) is exactly \(\tau\).  Thus \(q\) vanishes on \(\tau\) and
is strictly positive on every ray of \(\sigma\) not contained in
\(\tau\).  The principal toric open
\(D(\chi^q)\) contains \(O(\tau)\).  Its monoid is
\begin{equation}
\label{eq:p061-localized-toric-monoid}
 (\sigma^\vee\cap M)+\mathbf Z(-q)
   =\{m\in M: m|_\tau\geq0\}.
\end{equation}
For completeness, the inclusion from left to right is immediate.  For
the converse, evaluate \(m\) on the finitely many primitive ray
generators of \(\sigma\).  It is nonnegative on those in \(\tau\), while
\(q\) is positive on all the others.  Hence \(m+r q\in\sigma^\vee\) for
all sufficiently large integers \(r\), which proves
\eqref{eq:p061-localized-toric-monoid}.

Dualizing the chosen lattice splitting identifies the monoid on the
right with
\[
 (\tau^\vee\cap N_\tau^\vee)\oplus (N')^\vee.
\]
Consequently there is an algebraic isomorphism
\begin{equation}
\label{eq:p061-orbit-product-chart}
 D(\chi^q)\simeq
 X_{\tau,N_\tau}\times_K T_{N'},
\end{equation}
under which \(O(\tau)\) is the product of the torus-fixed point of
\(X_{\tau,N_\tau}\) and the split torus \(T_{N'}\).  Notice that
\(\tau\) is full-dimensional in \((N_\tau)_{\mathbf R}\), whether or not
it is simplicial.

Under \eqref{eq:p061-orbit-product-chart}, the point \(x\) has the form
\((a,y)\), with \(a\) the rational fixed point and
\(y\in T_{N'}^{\mathrm{an}}\).  The set
\(W\cap D(\chi^q)^{\mathrm{an}}\) is an ordinary open neighbourhood of
\((a,y)\).  Apply
\Cref{prop:p061-parameterized-toric-fixed-fibre}.  It factors \(h\), on
a product affinoid neighbourhood of \(x\), as the pullback of a
classically pluriharmonic function on a strict affinoid domain of the
split torus.  The torus is quasi-smooth, so the verified quasi-smooth
P06.1 theorem makes that function \(\mathbf R\)-PL.  Pullback stability
\cite[Lemma~5.4(1)]{GJR26} makes \(h\) \(\mathbf R\)-PL near \(x\).

Since \(x\) was arbitrary and \(\mathbf R\)-PL functions form a sheaf
\cite[\S5.2(4)]{GJR26}, \(h\) is \(\mathbf R\)-PL on \(W\).  Both \(h\)
and \(-h\) are classically psh.  On a good strict analytic space,
\cite[Theorem~4.3]{GR25Classical} identifies classical psh
\(\mathbf R\)-PL functions with semipositive ones.  Thus both signs are
semipositive, which is harmonicity in
\cite[Definition~7.2]{GJR26}.
\end{proof}

\begin{proof}[Proof of \Cref{cor:p061-split-toric-varieties}]
A split toric variety is covered by the affine toric opens attached to
the cones of its fan.  The first assertion therefore follows from
\Cref{thm:p061-split-toric-all-orbits} and sheafiness.  In the second
assertion, transfer a classically pluriharmonic function through the
given analytic local isomorphism and apply
\Cref{prop:p061-parameterized-toric-fixed-fibre}.  The quasi-smooth
P06.1 theorem applies to the \(Y\)-factor, and pullback through the local
isomorphism returns the claimed \(\mathbf R\)-PL and harmonic germ.
\end{proof}

\begin{proof}[Proof of \Cref{cor:p061-split-toric-quasismooth-products}]
Both assertions are local, so fix \((x,y)\in\Omega\) and first work on an
affine toric chart \(X_\sigma\) containing \(x\).  Let
\(\tau\preceq\sigma\) be the face whose orbit contains \(x\).  The
principal-open decomposition from the proof of
\Cref{thm:p061-split-toric-all-orbits} gives an ordinary neighbourhood of
\(x\) isomorphic to
\[
 X_{\tau,N_\tau}^{\mathrm{an}}\times_KT_{N'}^{\mathrm{an}},
\]
where \(\tau\) is full-dimensional and strongly convex in
\((N_\tau)_{\mathbf R}\), and where \(x=(a,t)\) with \(a\) the
torus-fixed point of \(X_{\tau,N_\tau}\).  Thus, after this restriction,
\(\Omega\) is an ordinary open neighbourhood of
\[
 (a,(t,y))\in
 X_{\tau,N_\tau}^{\mathrm{an}}
 \times_K\bigl(T_{N'}^{\mathrm{an}}\times_KY\bigr).
\]
The parameter \(T_{N'}^{\mathrm{an}}\times_KY\) is good, strict, and
quasi-smooth: a split torus is quasi-smooth, and quasi-smoothness is stable
under finite products.  Applying
\Cref{prop:p061-parameterized-toric-fixed-fibre} with this parameter factors
\(h\), on a product affinoid neighbourhood of \((x,y)\), as the pullback
of a classically pluriharmonic function on a quasi-smooth affinoid domain.
By \Cref{thm:p061-quasismooth-ph-rpl}, that function is
\(\mathbf R\)-PL.  Pullback stability
\cite[Lemma~5.4(1)]{GJR26} makes \(h\) \(\mathbf R\)-PL near \((x,y)\).

The same argument applies at every point.  Since \(\mathbf R\)-PL
functions form a sheaf \cite[\S5.2(4)]{GJR26}, \(h\) is
\(\mathbf R\)-PL on \(\Omega\).  Applying
\cite[Theorem~4.3]{GR25Classical} to both \(h\) and \(-h\) shows that
both are semipositive, which is Gubler--Jell--Rabinoff harmonicity by
\cite[Definition~7.2]{GJR26}.  Finally, transfer through an analytic local
isomorphism proves the stated local-model consequence.  This argument uses
an actual split toric product and a quasi-smooth parameter; it does not
assert the conclusion for nonsplit toric factors, non-quasi-smooth
parameters, or arbitrary nonproduct analytic germs.
\end{proof}

\begin{proof}[Proof of \Cref{cor:p061-toric-forms}]
Let \(\pi:X_L^{\mathrm{an}}\to X^{\mathrm{an}}\) be the scalar-extension
map.  Extension of scalars preserves classical pluriharmonicity in both
directions by \cite[Proposition~7.2(6)]{GR25Classical}.  Thus the pullback
of the given function is classically pluriharmonic on the inverse image
of the ordinary open subspace.  This inverse image is an ordinary open
subspace of the split toric variety \(X_L^{\mathrm{an}}\), so
\Cref{cor:p061-split-toric-varieties} makes the pullback
\(\mathbf R\)-PL.  Ground-field descent for \(\mathbf R\)-PL functions
\cite[Corollary~5.12]{GJR26} gives the assertion over \(K\).
Harmonicity follows, as above, from
\cite[Theorem~4.3]{GR25Classical} applied to both signs.
\end{proof}
